% Generated draft: requires mathematical and source review.
\documentclass[11pt]{article}
\usepackage[margin=25mm]{geometry}
\usepackage{fontspec}
\setmainfont{Latin Modern Roman}
\usepackage{amsmath,amssymb,mathtools,unicode-math}
\setmathfont{Latin Modern Math}
\usepackage{enumitem,needspace,xurl,hyperref}
\setlength{\parindent}{0pt}
\setlength{\parskip}{5pt}
\setlength{\emergencystretch}{3em}
\newcommand{\N}{\mathbb{N}}
\newcommand{\Z}{\mathbb{Z}}
\newcommand{\Q}{\mathbb{Q}}
\newcommand{\R}{\mathbb{R}}
\newcommand{\C}{\mathbb{C}}
\newcommand{\F}{\mathbb{F}}
\newcommand{\A}{\mathbb{A}}
\newcommand{\PP}{\mathbb{P}}
\newcommand{\E}{\mathbb{E}}
\newcommand{\eps}{\varepsilon}
\newcommand{\dd}{\,\mathrm{d}}
\begin{document}
\section*{UnsolvedMath ID: 2100604}
\textbf{Source problem code:} catalogue entry
\subsection{Definitions and mathematical statement}
An affine manifold has coordinate changes of the form $x\mapsto Ax+b$ with invertible linear part $A$; a foliation with affine leaves equips each leaf with such a structure. This does not by itself require compact leaves or tori. In the regular Liouville-integrable symplectic setting, a momentum map is a map
\[
F=(f_1,\ldots,f_n):(M^{2n},\omega)\longrightarrow\mathbb R^n
\]
with $\{f_i,f_j\}=0$ for all $i,j$ and $df_1,\ldots,df_n$ independent on the regular set. Its regular fiber components are Lagrangian, and compact connected regular fibers carry the affine-torus structure from action-angle theory. This is one conventional context, not a definition imposed on the source's unspecified moment map.

\textbf{Source question.} Which foliations with affine leaves can be described as the image of the moment map? A complete characterization requires specifying the ambient foliation, the class of moment maps, and the exact meaning of image. The image $F(M)$ lies in the target, whereas a foliation by fibers lies in the domain; these are different objects.

The supplied research summary reads the question as an inverse problem asking which affine-torus foliations arise as Liouville fibrations of integrable systems. Under that separate interpretation, one seeks necessary and sufficient conditions for a given foliation to be realized by connected fibers of an appropriate momentum map, with the prescribed affine structures. The original sentence neither says fibers nor restricts all affine leaves to compact tori. Both discrepancies remain unresolved, so this interpretation is not silently substituted for the literal question. The summary reports no verified complete characterization.
\subsection{Short English statement}
Characterize affine-leaf foliations realizable by the intended moment-map construction. The source says image, while its research summary discusses fiber foliations; this distinction must be resolved before a precise theorem is possible.
\subsection{Sources}
\begin{itemize}
\item \emph{Open Problems, Questions, and Challenges in Finite-Dimensional Integrable Systems} (2018), Problem 6.4, \url{https://arxiv.org/abs/1804.03737}; source-TeX extraction, accessed 2026-07-29. No author list is supplied in this record.
\item The supplied summary invokes Liouville torus fibrations, action-angle ideas, and the inverse problem for Lagrangian fibrations, but gives no particular characterization theorem or bibliographic reference.
\item Local registry \texttt{data/unsolvedmath/problems.json}, permanent ID $2100604$, source code catalogue entry, \emph{Open Problems in Integrable Systems --- Poisson geometry and action-angle variables}. The record labels this open and OPEN-TRIAGE, with NEEDS\_REVIEW source-status and rights qualifications.
\end{itemize}
\end{document}
